\documentclass[10pt]{article}
\usepackage[russian]{babel}
\usepackage[utf8]{inputenc}
\usepackage[T2A]{fontenc}
\usepackage{amsmath}
\usepackage{amsfonts}
\usepackage{amssymb}
\usepackage[version=4]{mhchem}
\usepackage{stmaryrd}
\usepackage{mathrsfs}
\usepackage{bbold}

\begin{document}
\section*{ГЛАВА 3}
\section*{АТОМАРНЫЕ ФУНКЦИИ}
Важную роль в теории приближений и численных методах играют алгебраические и тригонометрические (экспоненциальные) полиномы. Однако этот универсальный аппарат приближения имеет следующий (существенный для практики) недостаток - «нефинитность» классических полиномов. В численных же методах часто желательно применять «локальные» (т. е. финитные) функции с носителями малого диаметра, так как (см. §6 гл. 2) в этом случае в соответствующих матрицах большинство элементов равны нулю, матрица - «редкая» или даже «ленточная» и при том же объеме оперативной памяти ЭВМ можно использовать аппроксимирующие подпространства более высоких размерностей, а следовательно, получить более точные приближения. Известные и широко применяемые $[7,17,18,71,83,88,90]$ локальные функции-сплайны (например, $B$-сплайны Шенберга) не универсальны в том смысле, что для получения оптимального приближения функций большей гладкости нужны сплайны более высокой степени. Таким образом, классические алгебраические и тригонометрические полиномы универсальны, но не локальны, сплайны же локальны, но не универсальны. Естественно, возникает вопрос о построении пространств функций, которые были бы одновременно локальны и универсальны. Желательно к тому же, чтобы их можно было легко вычислять. Функции, о которых речь идет ниже, на наш взгляд, настолько просты и удобны в работе, что ими целесообразно пополнить класс «элементарных» функций.

Как же построены эти функции? Так как неуниверсальность сплайнов обусловлена уже их недостаточной гладкостью, то необходимо рассмотреть бесконечно дифференцируемые финитные функции. Финитные сплайны степени $n$, как хорошо известно, - это $n$ кратные свертки характеристических функций интервалов; эти финитные функции - бесконечнократные свертки характеристических функций интервалов. Для получения финитных функций длины интервалов должны достаточно быстро стремиться к нулю (сумма длин должна быть конечной). Если $\varphi(x)$ - искомая функция и\\
$\hat{\varphi}(t)$ - ee фурье-образ, то

$$
\hat{\varphi}(t)=\prod_{k=1}^{\infty} \frac{\sin \alpha_{k} t}{\alpha_{k} t} ;
$$

здесь $\alpha_{k}$ монотонно стремятся к нулю, $\alpha_{k}>0, \sum_{k=1}^{\infty} \alpha_{k}<\infty$. Теперь из всех монотонно стремящихся к нулю последовательностей $\alpha= =\left\{\alpha_{k}\right\}$ выберем те, для которых $\varphi(x)=\varphi_{\alpha}(x)$ обладает достаточно хорошими аппроксимационными свойствами. По-видимому, наиболее простой и удобной будет функция $\varphi(x)$, соответствующая геометрической прогрессии $\alpha_{k}=2^{-k}$. Такой выбор $\alpha_{k}$ по крайней мере естествен. Обозначим эту функцию up $(x)$. Таким образом,

$$
\operatorname{up}(x)=\frac{1}{2 \pi} \int_{-\infty}^{\infty} e^{-i t x} \prod_{k=1}^{\infty} \frac{\sin t 2^{-k}}{t 2^{-k}} d t .
$$

Первоначально функция up $(x)$ появилась при следующих обстоятельствах. В 1967 г. В. Л. Рвачев поставил такую задачу. Если $\varphi(x)$ - финитная дифференцируемая функция, имеющая один участок возрастания и один участок убывания («горб»), то ее производная $\varphi^{\prime}(x)$ состоит из «горба» и «ямы». Существует ли функция $\varphi(x)$, у которой «горб» и «яма» производной подобны «горбу» самой функции? На языке уравнений это означает: существует ли финитное решение уравнения $y^{\prime}(x)=a[y(2 x+1)-y(2 x-1)]$, в котором для определенности считаем, что носитель $\varphi(x)$ - отрезок $[-1,1]$ ?

В работе [52] доказаны существование и единственность такого финитного решения с интегралом, равным 1 , - это и есть функция up ( $x$ ). Таким образом, в то время как классические алгебраические и тригонометрические полиномы удовлетворяют однородным линейным дифференциальным уравнениям с постоянными коэффициентами, функция up ( $x$ ) (и другие аналогичные функции) удовлетворяет уравнениям вида

$$
L y(x)=\sum_{k=1}^{M} c_{k} y\left(a x-b_{k}\right),
$$

где $L$ - линейный дифференциальный оператор с постоянными коэффициентами. Эти уравнения близки к линейным дифференциальным уравнениям с постоянными коэффициентами в том смысле, что преобразование Фурье для них «работает» так же эффективно. Наблюдавшееся до сих пор некоторое невнимание к этим уравнениям и их решениям обусловлено, по-видимому, отсутствием их непосредственной физической интерпретации. В данной главе рассматриваются финитные решения (названные атомарными функциями) указанных выше уравнений с точки зрения их возможной применимости в конструктивной теории функций и численном анализе.

\section*{§ 1. ОПРЕДЕЛЕНИЕ АТОМАРНЫХ ФУНКЦИИ}
Атомарные функции - это финитные решения дифференциальнсфункциональных уравнений вида


\begin{equation*}
L y(x)=\lambda \sum_{k=1}^{M} c_{k} y\left(a x-b_{k}\right), \tag{3.1}
\end{equation*}


где $a>1, L$ - линейный обыкновенный дифференциальный оператор с постоянными коэффициентами.

Обозначим оператор

$$
y(x) \rightarrow \sum_{k=1}^{M} c_{k} y\left(a x-b_{k}\right)
$$

через $R$. Первый вопрос, который необходимо выяснить: какие условия нужно наложить на $L$ и $R$, чтобы уравнение (3.1) для некоторых $\lambda \neq 0$ имело ненулевое финитное решение?

Пусть $L=P(D)$, где $D=-i \frac{d}{d x}$, а

$$
P(\lambda)=\sum_{k=0}^{n} d_{k} \lambda^{k}
$$

Применим к обеим частям (3.1) преобразование Фурье

$$
\int_{-\infty}^{\infty} e^{i t x} L y(x) d x=\lambda \sum_{k=1}^{M} c_{k} \int_{-\infty}^{\infty} e^{i t x} y\left(a x-b_{k}\right) d x
$$

После очевидных выкладок получим


\begin{equation*}
P(t) F(t)=\lambda \sum_{k=1}^{M} \frac{c_{k}}{a} e^{i \frac{t b_{k}}{a}} F\left(\frac{t}{a}\right), \tag{3.2}
\end{equation*}


где

$$
F(t)=\int_{-\infty}^{\infty} e^{i t x} y(x) d x
$$

Положим

$$
\mathscr{E}(t)=\sum_{k=1}^{M} \frac{c_{k}}{a} e^{i \frac{t b_{k}}{a}}
$$

Рассмотрим сначала случай

$$
F(0)=\int_{-\infty}^{\infty} y(x) d x \neq 0
$$

Без ограничения общности можно считать, что


\begin{equation*}
F(0)=1 . \tag{3.3}
\end{equation*}


Тогда из (3.2) следует $P(0) F(0)=\lambda \mathscr{E}(0)$, т. е.


\begin{equation*}
d_{0}=\lambda \mathscr{E}(0) . \tag{3.4}
\end{equation*}


Пусть $P(\lambda)=d_{f} \lambda^{f}+\cdots+d_{n} \lambda^{n}$, где $d_{f} \neq 0\left(d_{0}=\cdots=d_{f-1}=\right.$ = 0). Тогда, последовательно дифференцируя (3.2) и полагая $t= =0$, получаем $f!d_{f}=\lambda \mathscr{E}(f)(0), \mathscr{E}(i)(0)=0,0 \leqslant j \leqslant f-1$. Отсюда


\begin{equation*}
\lambda=\frac{f!d_{f}}{g^{(f)}(0)} . \tag{3.5}
\end{equation*}


Пусть $A=\left\{\alpha_{1}, \ldots, \alpha_{n}\right\}$ - множество нулей многочлена $P(t)$ (с учетом кратности), $B=\left\{\beta_{1}, \ldots, \beta_{k}, \ldots\right\}$ - множество нулей функции $\mathscr{E}(t)$ (также с учетом кратности). Из $P(t) F(t)= =\lambda \mathscr{E}(t) F\left(t a^{-1}\right)$ следует


\begin{equation*}
P(t) F(t) P\left(t a^{-1}\right)=\lambda^{2} \mathscr{E}(t) \mathscr{E}\left(t a^{-1}\right) F\left(t a^{-2}\right) . \tag{3.6}
\end{equation*}


Так как (см. (3.2)) $P\left(t a^{-1}\right) F\left(t a^{-1}\right)=\lambda \mathscr{E}\left(t a^{-1}\right) F\left(t a^{-2}\right)$, то, умножив (3.2) на $P\left(t a^{-1}\right)$, получим (3.6)). Таким же образом получаем


\begin{equation*}
F(t) \prod_{k=0}^{n} P\left(t a^{-k}\right)=\lambda^{n+1} \prod_{k=0}^{n} \mathscr{E}\left(t a^{-k}\right) F\left(t a^{-n-1}\right) . \tag{3.7}
\end{equation*}


Поскольку $a>1$ и $F(0) \neq 0$, то начиная с некоторого $n$ имеем $F\left(t a^{-n-1}\right) \neq 0$ (для данного $t$ ). Отсюда каждый нуль функции $P\left(t a^{-m}\right)$ должен содержаться среди нулей функции $\mathscr{E}\left(t a^{-k(m)}\right)$ для некоторого $k(m)$.

Теорема 1. Необходимым и достаточным условием существования финитного ненулевого решения уравнения (3.1) $y(x)$ такого, что

$$
\int_{-\infty}^{\infty} y(x) d x=1
$$

является следующее условие (И): существует инъекция $i: A \rightarrow B$ такая, что для всякого $\alpha_{i} \in A$ существует натуральное $m$, для которого $i\left(\alpha_{j}\right)=a^{-m} \alpha_{j} \in$ В. Коэффициент $\lambda$ получается при этом из (3.5), а решение имеет вид

$$
y(x)=\frac{1}{2 \pi} \int_{-\infty}^{\infty} e^{-i x t} F(t) d t,
$$

где


\begin{equation*}
F(t)=\prod_{k=0}^{\infty} \frac{\lambda \xi\left(t a^{-k}\right)}{P\left(t a^{-k}\right)} . \tag{3.8}
\end{equation*}


Доказательство. Достаточность. Нужно показать, что $F(t)$ - целая функция экспоненциального типа первого порядка pocta: $|F(t)|<C e^{k|t|}$ и $F_{1}(t) \in L_{1}(-\infty, \infty),|F(t)|<C_{m} t^{-m}$, $t \in \mathbb{R}, m \in N$. Тогда $y(x)$ по теореме Винера - Пэли - финитная функция, и выполнение уравнения (3.1) проверяется непосредственно.

Пусть $|t| \leqslant R$. Тогда начиная с некоторого $N$ многочлен $P\left(t a^{-N}\right)$ в круге $|t| \leqslant R$ нулей, отличных от $t=0$, не имеет. Следовательно, $\frac{\lambda g\left(t a^{-k}\right)}{P\left(t a^{-k}\right)}$ при $k>N$ - аналитическая в круге $|t| \leqslant \leqslant R$ функция и

$$
\psi(t)=\prod_{k=N}^{\infty} \frac{\lambda \xi\left(t a^{-k}\right)}{P\left(t a^{-k}\right)}
$$

\begin{itemize}
  \item также целая функция [21, с. 11, теорема 2]. Это следует из того, что ряд
\end{itemize}

$$
\sum_{k=N}^{\infty}\left(\frac{\lambda \varepsilon\left(t a^{-k}\right)}{P\left(t a^{-k}\right)}-1\right)
$$

равномерно сходится в $|t| \leqslant R$, поскольку

$$
\left|\frac{\lambda E\left(t a^{-k}\right)}{P\left(t a^{-k}\right)}-1\right|<\frac{C_{p}}{a^{k}} .
$$

В силу условия (И) каждому нулю знаменателя в (3.8) соответствует свой нуль числителя, так что функция

$$
F(t)=\prod_{k=0}^{N-1} \frac{\lambda E\left(t a^{-k}\right)}{P\left(t a^{-k}\right)} \psi(t)
$$

также аналитическая в $|t| \leqslant R$ и, следовательно, целая, так как $R$ - любое.

Пусть $\mu=\max _{|t|<a^{2} R}|F(t)|$, где $R=\max \left|\alpha_{k}\right|$. Тогда

$$
\begin{aligned}
|F(t)| & \leqslant\left|\prod_{i=0}^{q_{0}} \frac{\lambda E\left(t a^{-j}\right)}{P\left(t a^{-j}\right)}\right| \mu ; \\
q_{0} & =\left[\log _{a} \frac{|t|}{a^{2} R}\right]+1 .
\end{aligned}
$$

Пусть $v=\min _{|t|>a R}|P(t)|>0$. Тогда при $|t|>a R$

$$
|F(t)|<\mu\left(\frac{|\lambda| \sum_{k=1}^{M}\left|\frac{c_{k}}{a}\right|}{v}\right)^{q_{0}} e^{|t| \frac{\max \left|b_{k}\right| a}{1-a}}<C e^{K|t|} .
$$

Рассмотрим поведение $F(t)$ на действительной оси. Пусть $\mu_{A}= =\max _{|t|<A}|F(t)|, t \in R, A \geqslant T ; T$ еыбирается из соотношения $|P(t)|> >\frac{d_{n}}{2}|t|^{n} \quad$ при $\quad|t| \geqslant T$ и $l=|\lambda| \sum_{k=1}^{M} \frac{\left|c_{k}\right|}{a}$. Тогда из

$$
F(t)=\frac{\lambda e(t)}{P(t)} F\left(\frac{t}{a}\right)
$$

следует

$$
|F(t)|<\frac{2 l}{d_{n} \mid t_{i}^{n}{ }^{n}} \mu_{A}
$$

при $A<|t|<a A, t \in \mathbb{R}$. Если

$$
A>\gamma>\sqrt[n]{\frac{\overline{2 l}}{d_{n}}},
$$

то $|F(t)|<\rho \mu_{A}$, где $0<\rho<1$. Следовательно, $\mu_{a A}=\mu_{A}$ и

$$
|F(t)|<\frac{2 l}{d_{n}|t|^{n}} \mu_{a}
$$

для $A<t<a^{2} A$. Таким образом,

$$
|F(t)|<\frac{2 l}{d_{n}|t|^{n}} \mu_{A}
$$

для всех $|t|>A, t \in \mathbb{R}$. Несложной модификацией этих рассуждений можно показать, что $F(t)$ на действительной оси убывает быстрей любой степени $t^{n}$. В самом деле, как уже показано ранее, на действительной оси $|F(t)|$ ограничена, а

$$
F(t)=\prod_{i=0}^{s} \frac{\lambda g\left(t a^{-i}\right)}{P\left(t a^{-j}\right)} F\left(t a^{-s-1}\right) .
$$

Таким образом, $y(x)$ - действительно финитная функция класса $C^{\infty}$. Достаточность доказана.

Необходимость. Из того, что $y(x)$ финитна, следует, что $F(t)$ целая. Поэтому должна существовать инъекция

$$
\gamma: \bigcup_{k=0}^{\infty} a^{k} A \rightarrow \bigcup_{k=0}^{\infty} a^{k} B
$$

(здесь $a^{k} A=\left\{a^{k} \alpha_{1}, \ldots, a^{k} \alpha_{n}\right\}$, аналогично понимается $a^{k} B$, объединение здесь и далее считается дизъюнктивным). Действительно, каждое $\mu \neq 0, \mu \in \mathbb{C}$, в множествах $\bigcup_{k=0}^{\infty} a^{k} A$ и $\bigcup_{k=0}^{\infty} a^{k} B$ может встречаться только конечное число раз, т. е.

$$
\begin{aligned}
& \mu \in \bigcup_{k=0}^{N_{1}} a^{k} A \quad \text { и } \quad \mu \notin \bigcup_{k=N_{1}+1}^{\infty} a^{k} A, \\
& \mu \in \bigcup_{k=1}^{N_{2}} a^{k} B \quad \text { и } \quad \mu \notin \bigcup_{k=N_{2}+1}^{\infty} a^{k} B .
\end{aligned}
$$

(Случай $\mu=0$ не вызывает затруднений, та́к как выше показано, что для $\mathscr{E}(t) \mu=0$ - нуль того же порядка, что и для $P(t)$ ). Тогда из (3.7) при $n>\max \left(N_{1}, N_{2}\right)$ следует, что $\mu$ встречается в $\bigcup_{k=0}^{\infty} a^{k} A$ не чаще, чем в $\bigcup_{k=0}^{\infty} a^{k} B$ (так как $F\left(\mu a^{-n-1}\right) \neq 0$ для достаточно больших n), и, следовательно, инъекция $\gamma$ существует.

Положим

$$
\begin{gathered}
S_{0}=\left\{\alpha_{i}: \alpha_{i} a^{-k} \notin A, k=1,2, \ldots\right\}, \\
S_{1}=\left\{\alpha_{i} \in A \backslash S_{0}: \alpha_{i} a^{-k} \notin A \backslash S_{0}, k=1,2, \ldots\right\}, \\
\dot{S}_{2}=\left\{\alpha_{i} \in A \backslash S_{0} \cup S_{1}: \alpha_{i} a^{-k} \in A \backslash S_{0} \cup S_{1}, k=1,2, \ldots\right\} ;
\end{gathered}
$$

$S_{l}$ - последнее непустое $S_{j}$. Необходимо по инъекции

$$
\gamma: \bigcup_{k=0}^{\infty} a^{k} A \rightarrow \bigcup_{k=0}^{\infty} a^{k} B
$$

построить инъекцию $i$. Начнем с $\alpha_{i} \in S_{0}$. Каждому из них инъективно ставим в соответствие $i\left(\alpha_{j}\right) \in B$, а именно: $\gamma\left(\alpha_{j}\right)=a^{m} \beta_{s}$. Положим $i\left(\alpha_{j}\right)=\rho_{s}$. Инъективность очевидна.

Теперь перестроим $\gamma$ в $\gamma_{1}$ следующим образом. Пусть $\alpha_{i_{r}} \in S_{0}$, $r=1, \ldots, r^{*}$. Расположим $j_{r}$ в некотором порядке. Рассмотрим $a^{k} \alpha_{i_{r}}, k=1,2, \ldots$ Положим для $\mu \in \bigcup_{k=0}^{\infty} a^{k} A$, где $\gamma_{1,0}=\gamma$,

$$
\gamma_{1, r}(\mu)= \begin{cases}\gamma_{1, r-1}(\mu), & \text { если } \quad \mu \neq a^{k} \alpha_{j_{r}} \text { и } \quad \gamma_{1, r-1}(\mu) \neq a^{l} i\left(\alpha_{j_{r}}\right) ; \\ a^{l} i\left(\alpha_{j_{r}}\right), & \text { если } \quad \mu=a^{l} \alpha_{j_{r}} ; \\ \gamma_{1, r-1}\left(a^{l} \alpha_{j_{r}}\right), & \text { если } \gamma_{1, r-1}(\mu)=a^{l} i\left(\alpha_{i_{r}}\right) .\end{cases}
$$

Пусть. $\gamma_{1}=\gamma_{1, r^{*}}$. Легко видеть, что для всех $\alpha_{i} \in S_{0} \gamma_{1}\left(a^{l} \alpha_{j}\right)= =a^{l} \gamma\left(\alpha_{j}\right)$.

Переходим к $S_{1}$. Полагаем $i(\alpha)$ для $\alpha \in S_{1}$ следующим образом. Пусть $\gamma_{1}(\alpha)=a^{m} \beta_{s}$. Тогда $i(\alpha)=\beta_{s}$. Перестроим инъекцию $\gamma_{1}$ в инъекцию $\gamma_{2}$ аналогично тому, как $\gamma$ перестраивалась в $\gamma_{1}$. В конце концов получим требуемую инъекцию $i$. Теорема доказана.

В теореме 1 рассматривался случай, когда $F(0)=1$. Пусть теперь $F(0)=0$. Тогда $F(t)=t^{r} G(t)$, где $G(0) \neq 0$. Без ограничения общности можно считать, что $G(0)=1$. Имеем

$$
\mathscr{P}(t) t^{r} G(t)=\lambda \mathscr{E}(t) \frac{t^{r}}{a^{r}} G\left(\frac{t}{a}\right)
$$

т. е.


\begin{equation*}
\mathscr{P}(t) G(t)=\frac{\lambda}{a^{r}} \mathscr{E}(t) G\left(\frac{t}{a}\right) \tag{3.9}
\end{equation*}


и, следовательно, приходим к условиям теоремы 1 с заменой $\lambda$ на $\lambda a^{-r}$. Поэтому

$$
\lambda=\frac{f!d_{f} a^{r}}{\mathscr{E}^{(f)}(0)}, \quad G(t)=\prod_{j=0}^{\infty} \frac{\lambda a^{-r} g_{\left(t a^{-i}\right)}}{\mathscr{P}\left(t a^{-i}\right)},
$$

и


\begin{equation*}
\varphi(x)=\frac{1}{2 \pi} \int_{-\infty}^{\infty} e^{-i t x} G(t) \mathrm{dt} \tag{3.10}
\end{equation*}


является решением уравнения


\begin{equation*}
L \varphi(x)=\lambda a^{-r} \sum_{k=1}^{M} \frac{c_{k}}{a} \varphi\left(a x-b_{k}\right) \tag{3.11}
\end{equation*}


Поскольку $F(t)=t^{r} G(t)$ - искомое решение уравнения (3.1), $y(x)$ имеет вид


\begin{equation*}
y(x)=C \varphi^{(k)}(x) . \tag{3.12}
\end{equation*}


Таким образом, приходим к следующей теореме.\\
Теорема 2. Необходимым и достаточным условием существования ненулевых финитных решений уравнения (3.1) является условие (И). При выполнении этого условия для

$$
\lambda_{l}=\frac{a^{l} d_{f} f^{\prime}}{g^{(f)}(0)}
$$

существуют финитные решения $C \varphi_{l}(x), l \in N$, где

$$
\varphi_{0}(x)=\frac{1}{2 \pi} \int_{-\infty}^{\infty} e^{-i t x} \prod_{j=0}^{\infty} \frac{\lambda_{0} g^{\circ}\left(t a^{-j}\right)}{\mathscr{P}\left(t a^{-j}\right)} d t, \quad \varphi_{l}(x)=\varphi_{0}^{(l)}(x) .
$$

Замечание. $\varphi_{l}(x)$ линейно независимы.

\section*{§ 2. УРАВНЕНИЕ ПЕРВОГО ПОРЯДКА}
Исследуем подробно случай, когда $L=P(D)$, где


\begin{gather*}
P(t)=-i t+\alpha,  \tag{3.13}\\
R: y(x) \rightarrow c_{1} y\left(a x-b_{1}\right)+c_{2} y\left(a x-b_{2}\right) . \tag{3.14}
\end{gather*}


Тогда

$$
\mathscr{E}(t)=\frac{c_{1}}{a} e^{i \frac{b_{1} t}{a}}+\frac{c_{2}}{a} e^{i \frac{b_{2} t}{a}} .
$$

Пусть $\alpha \neq 0$. Тогда $\lambda\left(\frac{c_{1}}{a}+\frac{c_{2}}{a}\right)=\alpha$ и


\begin{equation*}
\lambda=\frac{a \alpha}{c_{1}+c_{2}}, \quad A=\{i \alpha\} . \tag{3.15}
\end{equation*}


Исследуем нули $\mathscr{E}(t)$. Пусть $\mathscr{E}(t)=0$. Тогда


\begin{equation*}
e^{i \frac{\left(b_{1}-b_{2}\right) t}{a}}=-\frac{c_{2}}{c_{1}}, \quad \beta=\frac{\operatorname{Ln}\left(-\frac{c_{2}}{c_{1}}\right) a}{i\left(b_{1}-b_{2}\right)} . \tag{3.16}
\end{equation*}


Условие (И) означает, что существует $m \in N$ такое, что $i \alpha=\beta a^{n}$, т. е.


\begin{equation*}
\alpha=\frac{\operatorname{Ln}\left(-\frac{c_{2}}{c_{1}}\right) a^{m+1}}{b_{2}-b_{1}} . \tag{3.17}
\end{equation*}


Без ограничения общности можно считать $c_{1}=1$. Тогда


\begin{equation*}
c_{2}=-\exp \left[\left(b_{2}-b_{1}\right) \alpha a^{-m-1}\right] . \tag{3.18}
\end{equation*}


Таким образом, для данных $b_{1}, b_{2}$ и $\alpha \neq 0$ существует счетное число коэффициентов $c_{2, n}$, для которых (3.1), (3.13), (3.14) имеют ненулевые финитные решения.

Пусть $\alpha=0$. Тогда


\begin{gather*}
c_{1}+c_{2}=0 \\
\lambda\left(\frac{c_{1}}{a} \frac{i b_{1}}{a}+\frac{c_{2}}{a} \frac{i b_{2}}{a}\right)=-i \tag{3.19}
\end{gather*}


Положим для определенности $c_{1}=1$. При этом


\begin{equation*}
\lambda=\frac{a^{2}}{b_{1}-b_{2}} . \tag{3.20}
\end{equation*}


Условие (И) означает, что существует $\beta=0$. Нули функции $\mathscr{E}(t)$ даются формулой $\beta=\frac{\operatorname{Ln} 1 \cdot a}{i\left(b_{1}-b_{2}\right)}$, среди них 0 действительно содержится.

Для $\alpha=0$ получаем единственное решение с $\int_{-\infty}^{\infty} y(x) d x=1$, имеющее вид


\begin{equation*}
y(x)=\frac{1}{2 \pi} \int_{-\infty}^{\infty} e^{-i t x} \prod_{j=0}^{\infty} \frac{a}{b_{2}-b_{1}} \frac{e^{i b_{1} t a-j-1}-e^{i b_{2} t a-i-1}}{t a^{-j}} \mathrm{dt} . \tag{3.21}
\end{equation*}


Сделав замену $x_{1}=x-\frac{b_{1}+b_{2}}{2}$, уравнение (3.1) с условиями (3.13) и (3.14) всегда можно привести к виду

$$
y^{\prime}(x)+\alpha y(x)=\lambda\left(c_{1} y\left(a x-b_{1}\right)+c_{2} y\left(a x+b_{1}\right)\right) .
$$

Подставив $x_{2}=\mu x$, получим


\begin{equation*}
y^{\prime}(x)+\alpha y(x)=\lambda(y(a x-1)+c y(a x+1)) . \tag{3.22}
\end{equation*}


Для $\alpha=0$ уравнение (3.22) имеет вид


\begin{equation*}
y^{\prime}(x)=\lambda(y(a x-1)-y(a x+1)) . \tag{3.23}
\end{equation*}


При этом


\begin{equation*}
y(x)=\frac{1}{2 \pi} \int_{-\infty}^{\infty} e^{i t x} \prod_{i=1}^{\infty} \frac{\sin t a^{-j}}{t a^{-j}} \mathrm{dt} ; \quad \lambda=-\frac{a^{2}}{2} . \tag{3.24}
\end{equation*}


При $a=2$ получаем $y(x)=u p(x)$. Обозначим функцию, определяемую формулой (3.24), $h_{a}(x)$. Выяснить вид носителя функции $h_{a}(x)$ можно, например, следующим образом. Как известно, функция $\frac{\sin t}{t}$ является характеристической функцией случайной величины, равномерно распределенной (р. р.) на отрезке [-1,1]. Тогда функция


\begin{equation*}
F_{a}(t)=\prod_{j=1}^{\infty} \frac{\sin t a^{-i}}{t a^{j}} \tag{3.25}
\end{equation*}


\begin{itemize}
  \item характеристическая функция случайной величины
\end{itemize}


\begin{equation*}
\xi(a)=\sum_{j=1}^{\infty} a^{-i} \xi_{j}, \tag{3.26}
\end{equation*}


где $\left\{\xi_{j}\right\}$ - последовательность независимых р. р. на $[-1,1]$ случайных величин. Функция $h_{a}(x)$ - плотность случайной величины $\xi(a)$. Теперь очевидно, что


\begin{equation*}
\operatorname{supp} h_{a}(x)=[-b, b], \tag{3.27}
\end{equation*}


где $b=\sum_{j=1}^{\infty} a^{-j}=\frac{1}{a-1}$. Тогда


\begin{equation*}
F_{a}(t)=\frac{\sin t a^{-1}}{t a^{-1}} F\left(t a^{-1}\right) . \tag{3.28}
\end{equation*}


Пусть $F_{a}(t)=\sum_{k=0}^{\infty} c_{k}(a) t^{k}$. Из (3.25) следует, что $c_{2 k+1}(a)=0$, так что


\begin{equation*}
F_{a}(t)=\sum_{k=0}^{\infty} c_{2 k}(a) t^{2 k} \tag{3.29}
\end{equation*}


Поскольку


\begin{equation*}
\frac{\sin t a^{-1}}{t a^{-1}}=\sum_{k=0}^{\infty}(-1)^{k} \frac{t^{2 k} a^{-2 k}}{(2 k+1)!}, \tag{3.30}
\end{equation*}


то (см. (3.28))


\begin{equation*}
c_{2 k}(a)=\sum_{j=0}^{k} \frac{c_{2 k-2 j}(-1)^{j} a^{-2 j}}{a^{2 k-2 j}(2 j+1)!}, \tag{3.31}
\end{equation*}


откуда

$$
c_{2 k}(a)\left(1-a^{-2 k}\right)=\sum_{j=1}^{k} \frac{c_{2 k-2 j}(a)(-1)^{j}}{(2 j+1)!} a^{-2 k}
$$

и, наконец,


\begin{equation*}
c_{2 k}(a)=\frac{1}{a^{2 k}-1} \sum_{l=0}^{k-1} \frac{c_{2 l}(a)(-1)^{k-l}}{(2 k-2 l+1)!} . \tag{3.32}
\end{equation*}


Нетрудно заметить, что $c_{0}(a)=1$. Из (3.32) следует

$$
\left|c_{2 k}(a)\right| \leqslant \frac{1}{a^{2 k}-1} \sum_{l=0}^{k-1} \frac{1}{(2 k-2 l+1) \mid} \max _{i<k}\left|c_{2 j}(a)\right| .
$$

Поскольку

$$
\sum_{l=0}^{k-1} \frac{1}{(2 k-2 l+1)!}<\sum_{j=0}^{\infty} \frac{1}{(2 j+1)!}=C_{1}<\infty,
$$

то начиная с некоторого $k \quad c_{2 k}(a)<\max _{i<k}\left|c_{2 j}(a)\right|$, откуда $\left|c_{2 k}(a)\right|< <C$ для всех $k$. Поэтому

$$
\left|c_{2 k}(a)\right| \leqslant\left(\prod_{j=1}^{k} \frac{1}{a^{2 i}-1}\right) C_{1}^{k} C_{2}
$$

и, следовательно, $\left|c_{2 k}(a)\right| \rightarrow 0$ при $k \rightarrow \infty$. Отсюда и из (3.32) следует


\begin{equation*}
\left|c_{2 k}(a)\right| \leqslant C_{3} \prod_{j=1}^{k} \frac{1}{\left(a^{2 j}-1\right)} \tag{3.33}
\end{equation*}


так как

$$
v_{k}=\sum_{l=0}^{k-1} \frac{\left|c_{2 l}(a)\right|}{(2 k-2 l+1) \mid} \rightarrow 0
$$

при $k \rightarrow \infty$, и начиная с некоторого $k v_{k}<1$. Из (3.33) нетрудно получить


\begin{equation*}
\left|c_{2 k}(a)\right|<c_{4} a^{-\sum_{i=1}^{k} 2 j}=C_{4} a^{-(k+1) k} \tag{3.34}
\end{equation*}


Оценку (3.34) можно еще улучшить. Скорость убывания $c_{2 k}(a)$ в дальнейшем понадобится, например, для оценки остаточного члена при разложении функции up ( $x$ ) в ряд специального вйда.

Уточним поведение $v_{k}$ при $k \rightarrow \infty$. Из (3.34) получаем

$$
\begin{gathered}
v_{k}<\sum_{l=0}^{k-1} \frac{C_{4} a^{-l(l+1)}}{(2 k-2 l+1)!}=C_{4} \sum_{j=1}^{k} \frac{a^{-(k-i+1)(k-j)}}{(2 j+1)!}= \\
=C_{4}\left(\sum_{i=1}^{\left[\frac{k}{2}\right]} \frac{a^{-(k-i+1)(k-j)}}{(2 j+1)!}+\sum_{i=\left[\frac{k}{2}\right]+1} \frac{a^{-(k-i+1)(k-j)}}{(2 j+1)!}\right) \leqslant \\
\leqslant C_{4}\left(a^{-\frac{k^{2}}{4}} \sum_{j=1}^{\left.\frac{k}{2}\right]} \frac{1}{(2 j+1)!}+\sum_{i=\left[\frac{k}{2}\right]+1}^{k} \frac{1}{(2 j+1)!}\right) \leqslant \\
\leqslant C_{4}\left(C_{5} a^{-\frac{k^{2}}{4}}+C_{6} \frac{1}{(k+1)!}\right) \leqslant C_{7} \frac{1}{(k+1)!} .
\end{gathered}
$$

Поскольку (см. (3.32)) $\left|c_{2 k}(a)\right| \leqslant \frac{1}{a^{2 k}-1} v_{k}$, то


\begin{equation*}
\left|c_{2 k}(a)\right| \leqslant \prod_{j=1}^{k} \frac{!}{a^{2 j}-1} \frac{C_{7}^{k}}{\prod_{j=1}^{k}(j+1)!} \leqslant C_{8} a^{-k(k+1)}\left(\prod_{j=1}^{k} j!\right)^{-1} . \tag{3.35}
\end{equation*}


Так как


\begin{equation*}
c_{2 k}(a)=\frac{F_{a}^{(2 k)}(0)}{(2 k)!}, \tag{3.36}
\end{equation*}


a


\begin{equation*}
F_{a}^{(2 k)}(0)=(-1)^{k} \int_{-\infty}^{\infty} x^{2 k} h_{a}(x) \mathrm{dx} \tag{3.37}
\end{equation*}


то


\begin{equation*}
\left|\int_{-\infty}^{\infty} x^{2 k} h_{a}(x) \mathrm{dx}\right| \leqslant C_{8} a^{-k(k+1)}(2 k)!\left(\prod_{j=1}^{k} j!\right)^{-1} \leqslant C_{9} a^{-k(k+1)} . \tag{3.38}
\end{equation*}


Рассмотрим $h_{a}(x)$ при $a>2$. Из (3.23) и (3.27) следует $h_{a}^{\prime}(x) \equiv \equiv 0$ на $\left[-b_{1}, b_{1}\right]$, где $b_{1}=b-\frac{2 b}{a}=b\left(1-\frac{2}{a}\right)>0$, и, значит, $h_{a}(x) \equiv$ const на $\left[-b_{1}, b_{1}\right]$. Аналогично получим, что на двух интервалах общей длины $2 \times 2 b_{1} a^{-1}$ функция $h_{a}(x)$ - многочлен первой степени, на четырех интервалах общей длины $4 \times 2 b_{1} a^{-2} h_{a}(x)$ - многочлен второй степени и т. д. Мера множества $P o l_{a}$ тех $x \in \operatorname{supp} h_{a}(x)$, у которых есть окрестность, где $h_{a}(x)$ - многочлен, составляет:


\begin{equation*}
\mu\left(P o l_{a}\right)=2 b_{1}+2 b_{1} \frac{2}{a}+2 b_{1}\left(\frac{2}{a}\right)^{2}+\cdots=2 b_{1} \frac{1}{1-\frac{2}{a}}=2 b, \tag{3.39}
\end{equation*}


т. е. $\mu\left(\right.$ Pol $\left._{a}\right)=\mu\left(\operatorname{supp} h_{a}(x)\right)$.

Таким образом, $h_{a}(x)$ при $a>2$ на множестве полной меры многочлен, а на оставшемся множестве sing ${ }_{a}$ (нигде не плотном, типа известного канторова множества) меры 0 - неаналитическая функция (действительно, ряд Тейлора функции $h_{a}(x)$ в точках множества $\operatorname{sing}_{a}$ либо состоит из конечного числа членов и не сходится, следовательно, к $h_{a_{s}}^{\mathrm{f}}(x)$, либо имеет нулевой радиус сходимости). Итак, в полученном примере функций класса $C^{\infty}(\mathbb{R})$ почти в каждой точке ряд Тейлора - многочлен, но функции тем не менее не являются многочленами и не аналитичны.

Поскольку кусочно-полиномиальные функции называются сплайнами, то функции $h_{a}(x)$ при $a>2$ естественно причислить к сплайнам класса $C^{\infty}$.

Вычислим значения $h_{a}(x)$ при $a>2$. Пусть Sing ${ }_{a}^{1}$ - множество тех точек из Sing ${ }_{a}$, в которых ряд Тейлора - многочлен. Поскольку $h_{a}(x)$ - четная функция, рассмотрим только интервал $[-b, 0]$. Из (3.22) получаем

$$
\begin{aligned}
h\left(-b_{1}\right) & =\int_{-b}^{b_{1}} h_{a}^{\prime}(x) \mathrm{dx}=-\lambda \int_{-b}^{b_{1}} h_{a}(a x+1) \mathrm{dx}= \\
& =-\frac{\lambda}{a} \int_{-b}^{b} h_{a}(x) \mathrm{dx}=-\frac{\lambda}{a}
\end{aligned}
$$

Таким образом, на $\left[-b_{1}, b_{1}\right] \quad h_{a}(x)=-\frac{\lambda}{a}$, а так как (см. (3.20)) $\lambda=-\frac{a^{2}}{2}$, то


\begin{equation*}
h_{a}(x)=\frac{a}{2} . \tag{3.40}
\end{equation*}


Поэтому


\begin{equation*}
\int_{-b_{1}}^{b_{1}} h_{a}(x) \mathrm{dx}=a b_{1}=\frac{a-2}{a-1} \tag{3.41}
\end{equation*}


и, следовательно,


\begin{equation*}
\int_{-b}^{b_{1}} h_{a}(x) \mathrm{dx}=\frac{1}{2}-\frac{1}{2} \frac{a-2}{a-1} . \tag{3.42}
\end{equation*}


Тогда


\begin{gather*}
h_{a}\left(-b+\frac{b-b_{1}}{a}\right)=\int_{-b}^{-b+\frac{b-b_{1}}{a}} h_{a}^{\prime}(x) \mathrm{dx}=-\lambda \int_{-b}^{-b+\frac{b-b_{1}}{a}} h_{a}(a x+1) \mathrm{dx}= \\
=-\frac{\lambda}{a} \int_{-b}^{b_{1}} h_{a}(x) \mathrm{dx}=\frac{a}{2}\left(\frac{1}{2}-\frac{1}{2} \frac{a-2}{a-1}\right) . \tag{3.43}
\end{gather*}


Далее,


\begin{gather*}
h_{a}^{\prime}\left(-b+\frac{b-b_{1}}{a}\right)=\int_{-b}^{-b+\frac{b-b_{1}}{a}} h_{a}^{\prime \prime}(x) \mathrm{d} x= \\
=\lambda^{2} a \int_{-b}^{-b+\left(b-b_{1}\right) a^{-1}} h_{a}\left(a^{2} x+a+1\right) \mathrm{d} x=\frac{\lambda^{2}}{a} \int_{-b}^{b} h_{a}(x) \mathrm{d} x=\frac{a^{3}}{4} . \tag{3.44}
\end{gather*}


Поэтому на $\left[-b+\frac{b-b_{1}}{a},-b_{1}-\frac{b-b_{1}}{a}\right]$.


\begin{equation*}
h_{a}(x)=\frac{a^{3}}{4}\left(x+b-\frac{b-b_{1}}{a}\right)+\frac{a}{2}\left(\frac{1}{2}-\frac{1}{2} \frac{a-2}{a-1}\right) . \tag{3.45}
\end{equation*}


Продолжая таким же образом, можно вычислить коэффициенты всех многочленов, из которых «составлена» функция $h_{a}(x)$, причем эти коэффициенты будут рациональным образом выражаться через $a$.

Между моментами функции $h_{a}(x) \int_{-b}^{b} x^{2 n} h_{a}(x) \mathrm{dx}$ (а следовательно, и коэффициентами $c_{2 k}(a)$ ) и значениями функции $h_{a}(x)$ в некоторых точках существует следующая связь. По формуле Тейлора

$$
h_{a}(x)=\sum_{k=0}^{n} \frac{h_{a}^{(k)}(-b)}{k!}(x+b)^{k}+R_{n}(x),
$$

где


\begin{equation*}
R_{n}(x)=\frac{1}{n!} \int_{-b}^{x} h_{a}^{(n+1)}(\tau)(x-\tau)^{n} \mathrm{~d} \tau \tag{3.46}
\end{equation*}


Но $h_{a}^{(k)}(-b)=0, k \in N$. Поэтому


\begin{equation*}
h_{a}(x)=R_{n}(x) \tag{3.47}
\end{equation*}


Из (3.46) и (3.23) видно, что значения $h_{a}(x)$ в точках $x \in \operatorname{Sing}_{a}^{1}$ есть линейные комбинации моментов функции $h_{a}(x)$.

Найдем выражения для $h_{a}^{(n)}(x)$ и опишем точки $x \in \operatorname{Sing}_{a}^{1}$. Из (3.23) следует


\begin{align*}
h_{a}^{\prime \prime}(x) & =\lambda^{2} a\left(h_{a}(a x+a+1)-h_{a}\left(a^{2} x+a-1\right)\right)- \\
& \left.-h_{a}\left(a^{2} x-a+1\right)+h_{a}\left(a^{2} x-a-1\right)\right) . \tag{3.48}
\end{align*}


В итоге, продолжая дифференцировать $h_{a}(x)$ с использованием (3.23), получаем


\begin{equation*}
h_{a}^{(n)}(x)=|\lambda|^{n} a^{\frac{n(n-1)}{2}} \sum_{k=1}^{2^{n}} \delta_{k} h_{a}\left(a^{n} x+\sum_{j=1}^{n} a^{j-1}(-1)^{p_{j}(k-1)}\right), \tag{3.49}
\end{equation*}


где $p_{j}(k)$ - число, состоящее в $j$-м разряде двоичного разложения числа $k$ (т. е. числа $k$, записанного в двоичной системе счисления). Числа $\delta_{k}$ определяются следующим рекуррентным соотношением: $\delta_{2 k}=-\delta_{k}, \delta_{2 k-1}=\delta_{k}, \delta_{1}=1$.

Теперь опишем точки множества Sing ${ }_{a}^{1}$. Точка $x_{0} \in \operatorname{Sing}_{a}^{1}$ в том и только в том случае, если существует натуральное $n$ и набор знаков + или - такой, что

$$
a^{n} x_{0}+\sum_{k=0}^{n-1}\left( \pm a^{k}\right)= \pm b
$$

т. е.


\begin{equation*}
x_{0}=\left( \pm b+\sum_{k=0}^{n-1}\left( \pm a^{k}\right)\right) a^{-n} \tag{3.50}
\end{equation*}


Пусть $x_{0} \in$ Sing $_{a}^{1}$. Величины $h_{a}^{(l)}\left(x_{0}\right)$ согласно (3.49) выражаются через значения $h_{a}(x)$ в точках вида

$$
x_{1}=a^{l} x_{0}+\sum_{i=0}^{l-1}( \pm a)
$$

Подставим вместо $x_{0}$ его значение из (3.50):

$$
\begin{aligned}
& x_{1}=a^{l-n}\left( \pm b+\sum_{k=0}^{n-1}\left( \pm a^{k}\right)\right)+\sum_{i=0}^{l-1}\left( \pm a^{i}\right)= \\
& =a^{l-n}\left( \pm b+\sum_{k=0}^{n-1}\left( \pm a^{k}\right)+\sum_{k=n-l}^{n-1}\left( \pm a^{k}\right)\right)
\end{aligned}
$$

Достаточно считать, что $0<l<n$, поскольку при $n \leqslant l h_{a}^{(l)}\left(x_{0}\right)= =0$. Покажем, что для тех наборов знаков в сумме $\Sigma_{\mathrm{I}}=\sum_{k=0}^{n-1}\left( \pm a^{k}\right)$, для которых при $n-l \leqslant k<n-1$ хотя бы один из знаков при той же степени $a$ в сумме

$$
\sum_{\mathrm{II}}=\sum_{k=n-l}^{n-1}\left( \pm a^{k}\right)
$$

не противоположен, выполняется неравенство $\left|x_{1}\right| \geqslant b$ и, следовательно, $h_{a}\left(x_{1}\right)=0$. Действительно, пусть $k_{1}, k_{1} \geqslant n-l$, наибольшее из тех $k$, для которых у данных наборов знаков в $\Sigma_{\text {I }}$ и $\Sigma_{\text {II }}$ знаки при $a^{k}$ одинаковы. Тогда

$$
\begin{aligned}
\left|\Sigma_{1}+\Sigma_{11}\right|> & 2 a^{k_{1}}-2 \sum_{k=n-l}^{k_{1}} a^{k}-\sum_{k=0}^{n-l-1} a^{k}=\frac{2\left(a^{k_{1}+1}-2 a^{k_{1}}\right)}{a-1}+ \\
& +\frac{a^{n-l}}{a-1}+\frac{1}{a-1} \geqslant \frac{a^{n-l}}{a-1}+b
\end{aligned}
$$

(так как $a>2$ ). Поэтому

$$
\left|a^{l-n}\left( \pm b+\sum_{\mathrm{I}}+\sum_{\mathrm{II}}\right)\right|>\frac{1}{a-1}=b .
$$

Для того (единственного) набора знаков в $\Sigma_{\text {II }}$, у которого знаки при $k \geqslant n-l$ противоположны знакам $\Sigma_{\mathrm{I}}$,

$$
x_{1}=a^{l-n}\left( \pm b+\sum_{k=0}^{n-l-1}( \pm a)\right)
$$

и в силу (3.50) $x_{1} \in \mathrm{Sing}_{a}^{1}$. Следовательно, и производные функции $h_{a}(x)$ в точках Sing ${ }_{a}^{1}$ рациональным образом выражаются через $a$ и $c_{2 k}(a)$ (т. е. в конце концов через $a$ ).

Пусть $1 \leqslant m \leqslant 2^{n}$, а $x_{m}^{(n)} \in \operatorname{Sing}_{a}^{1}$ имеет вид


\begin{equation*}
x_{m}^{(n)}=\left(b-\sum_{k=0}^{n-1} a^{k}(-1)^{p_{k+1^{(m-1)}}}\right) a^{-n} . \tag{3.51}
\end{equation*}


Тогда согласно (3.47) и (3.49)


\begin{gather*}
h_{a}\left(x_{m}^{(n)}\right)=\frac{1}{(n-1)!} \int_{-b}^{x_{m}^{(n)}} h_{a}^{(n)}(\tau)\left(x_{m}^{(n)}-\tau\right)^{n-1} \mathrm{~d} \tau= \\
=\frac{|\lambda|^{n} a^{\frac{n(n-1)}{2}}}{(n-1)!} \int_{-b}^{x_{m}^{(n)}} \sum_{k=1}^{m} \delta_{k} h_{a}\left(a^{n} \tau+\sum_{j=1}^{n} a^{j-1}(-1)^{p_{j}(k-1)}\right) \times \\
\times\left(x_{m}^{(n)}-\tau\right)^{n-1} d \tau=\frac{|\lambda|^{n} a^{\frac{(n-1) n}{2}}}{(n-1)!} \sum_{k=1}^{m} \delta_{k} \int_{-b}^{x_{m}^{(n)}} h_{a}\left(a^{n} \tau+\right. \\
\left.+\sum_{j=1}^{n} a^{j-1}(-1)^{p_{j}(k-1)}\right)\left(x_{m}^{(n)}-\tau\right)^{n-1} \mathrm{~d} \tau . \tag{3.52}
\end{gather*}


Сделаем в $k$-м интеграле в (3.52) замену

$$
t=a^{n} \tau+\sum_{i=1}^{n} a^{i-1}(-1)^{p_{j}(k-1)}
$$

Получим


\begin{align*}
h_{a}\left(x_{m}^{(n)}\right)= & \frac{|\lambda|^{n} a^{\frac{n(n-1)}{2}}}{(n-1)!} a^{-n} \sum_{k=1}^{m} \delta_{k} \int_{-b}^{b} h_{a}(t)\left(x_{m}^{(n)}-(t+\right. \\
& \left.\left.+\sum_{i=0}^{n-1} a^{l}(-1)^{p_{i+1}(k-1)}\right) a^{-n}\right)^{n-1} \mathrm{dt} . \tag{3.53}
\end{align*}


Отсюда $\left(|\lambda|=\frac{a^{2}}{2}\right)$\\
$h_{a}\left(x_{m}^{(n)}\right)=\frac{a^{2 n} a^{\frac{n(n-1)}{2}}}{(n-1)!2^{n}} a^{-n} a^{-n(n-1)} \sum_{k=1}^{m} \delta_{k} \sum_{l=0}^{\left[\frac{n-1}{2}\right]} C_{n-1}^{2 l} \beta_{m, k, n}^{n-1-2 l} \int_{-b}^{b} h_{a}(t) t^{2 l} \mathrm{dt}=$\\
$=\frac{1}{(n-1)!2^{n}} a^{-\frac{n(n-3)}{2}\left[\frac{n-1}{2}\right]} \sum_{l=0}^{2 l} C_{n-1}^{2 l}(-1)^{l}(2 l)!c_{2 l}(a) \sum_{k=1}^{m} \delta_{k} \beta_{m, k, n}^{n-1-2 l}$,\\
где

$$
\beta_{m, k, n}=b+\sum_{j=0}^{n-1} a^{i}\left((-1)^{p_{j+1}(k-1)}-(-1)^{p_{i+1}(m-1)}\right) .
$$

Например, для $m=1$, т. е. $x_{1}^{(n)}=-b+2 b a^{-n}$,


\begin{equation*}
h_{a}\left(x_{1}^{(n)}\right)=\frac{a^{-\frac{n(n-3)}{2}}}{(n-1)!2^{n}}\left[\frac{n-1}{2}\right] \sum_{l=0}^{2 l} C_{n-1}(-1)^{l}(2 l)!c_{2 l}(a) b^{n-1-2 l} . \tag{3.55}
\end{equation*}


Выведем формулу для $h_{a}^{(l)}\left(x_{m}^{(n)}\right)$. Из (3.49) следует


\begin{equation*}
h_{a}^{(l)}\left(x_{m}^{(n)}\right)=|\lambda|^{l} a^{\frac{l(l-1)}{2}} \delta_{s} h_{a}\left(a^{l} x_{m}^{(n)}+\sum_{j=0}^{l-1} a^{l}(-1)^{p_{i+1}(s-1)}\right), \tag{3.56}
\end{equation*}


где $s$ определяется из формулы $p_{1}(s-1)=p_{i+n-l}(m-1)$ : $s=\left[m 2^{l-n}\right]$. Таким образом,


\begin{gather*}
h_{a}^{(l)}\left(x_{m}^{(n)}\right)=|\lambda|^{l^{\frac{l(l-1)}{2}}} \delta_{s} h_{a}\left(a ^ { l - n } \left(b-\sum_{k=0}^{n-1} a^{k}(-1)^{p_{k+1}(m-1)}+\right.\right. \\
\left.\left.+\sum_{j=0}^{l-1} a^{l}(-1)^{p_{i+1}(s-1)}\right)\right)=\frac{a^{\frac{l(l+3)}{2}}}{2^{l}} \delta_{s} h_{a}\left(x_{r}^{(n-l)}\right), \tag{3.57}
\end{gather*}


где $r=m-\left[m 2^{l-n}\right] 2^{n-l}+1$. Окончательно получаем


\begin{gather*}
h_{a}^{(l)}\left(x_{m}^{(n)}\right)=\frac{a^{-\frac{(n-l)(n-l+3)}{2}+\frac{l(l+3)}{2}}}{2^{n}(n-l-1)!} \delta_{s} \times \\
\times \sum_{j=0}^{\left[\frac{n-l-1}{2}\right]} C_{n-l-1}^{2!}(-1)^{j}(2 j)!c_{2 j}(a) \sum_{k=1}^{r} \delta_{k} \beta_{r, k, n-l}^{n-l-1-2 j} . \tag{3.58}
\end{gather*}


Итак, на интервале $\left[x_{m}^{(n)}, x_{m}^{(n)}+2 b_{1} a^{-n+1}\right]$ функция $h_{a}(x)$ - многочлен степени $n-1$ вида


\begin{equation*}
h_{a}(x)=\sum_{l=0}^{n-1} \frac{h_{a}^{(l)}\left(x_{m}^{(n)}\right)}{l l}\left(x-x_{m}^{(n)}\right)^{l}, \tag{3.59}
\end{equation*}


где $h_{a}^{(l)}\left(x_{m}^{(n)}\right)$ вычисляются по формуле (3.58).

\section*{§ 3. ФУНКция up( $x$ )}
Рассмотрим $h_{a}(x)$ при $a=2$. По определению $h_{2}(x)=\operatorname{up}(x)$. Функция up $(x)$, введенная в работе [52], исследовалась в работах [53-57, 59, 60, 66]. Из (3.24) следует


\begin{equation*}
\operatorname{up}(x)=\frac{1}{2 \pi} \int_{-\infty}^{\infty} e^{i t x} \prod_{k=1}^{\infty} \frac{\sin t 2^{-k}}{t 2^{-k}} \mathrm{dt} . \tag{3.60}
\end{equation*}


Функция up $(x)$ - решение уравнения


\begin{equation*}
y^{\prime}(x)=2 y(2 x+1)-2 y(2 x-1), \tag{3.61}
\end{equation*}


причем


\begin{equation*}
\operatorname{supp} \operatorname{up}(x)=[-1,1] . \tag{3.62}
\end{equation*}


Производная $n$-го порядка функции up $(x)$ вычисляется по формуле


\begin{equation*}
\operatorname{up}^{(n)}(x)=2^{C_{n+1}^{2}} \sum_{k=1}^{2^{n}} \delta_{k} \operatorname{up}\left(2^{n} x+2^{n}+1-2 k\right) . \tag{3.63}
\end{equation*}


Легко заметить, что


\begin{equation*}
\left\|\operatorname{up}^{(n)}(x)\right\|_{C[-1,1]}=\left\|\operatorname{up}^{(n)}(x)\right\|_{L_{1}[-1,1]}=2^{C_{n+1}^{2}} . \tag{3.64}
\end{equation*}


В точках вида $k 2^{-n}$ (т. е. двоично-рациональных) ряд Тейлора функции up $(x)$ - многочлен степени $n$. Нетрудно убедиться, что в остальных точках носителя up $(x)$ - отрезка $[-1,1]$ - ряд Тейлора этой функции имеет нулевой радиус сходимости. Таким образом, в то время как функции $h_{a}(x)$ при $a>2$ аналитичны на множестве полной меры, функция up ( $x$ ) неаналитична ни в одной точке носителя. Повторяя рассуждения § 2, можно вычислить значения функции ир ( $x$ ) в двоично-рациональных точках, выразив их через момен-

ты функции ир ( $x$ ). Приведем формулы для этих моментов:


\begin{gather*}
2 b_{2 k}=a_{2 k}=\int_{-1}^{1} x^{2 k} \text { up }(x) \mathrm{dx}=(-1)^{k} c_{2 k}(2)(2 k)!  \tag{3.65}\\
b_{2 k+1}=\int_{0}^{1} x^{2 k+2} \text { up }(x) \mathrm{dx}=-\int_{0}^{1} \frac{x^{2 k+2}}{2 k+2} \text { up }^{\prime}(x) \mathrm{dx}= \\
=\frac{1}{k+1} \int_{0}^{1} x^{2 k+2} \text { up }(2 x-1) \mathrm{dx}=\frac{1}{2 k+2} \int_{-1}^{1}\left(\frac{u+1}{2}\right)^{2 k+2} \text { up }(u) d u= \\
=\frac{1}{2^{2 k+2}(2 k+2)} \sum_{j=0}^{k+1} C_{2 k+2}^{2 j} a_{2 j} \tag{3.66}
\end{gather*}


Теперь по формуле Тейлора получаем


\begin{gather*}
\operatorname{up}\left(-1+k 2^{-n}\right)=\frac{1}{n!} \int_{-1}^{-1+k 2^{-n}} \operatorname{up}^{(n+1)}(\tau)\left(-1+k 2^{-n}-\tau\right)^{n} d \tau= \\
=\frac{2^{C_{n+2}^{2}}}{n!} \sum_{j=1}^{k} \delta_{j}^{-1+k 2^{-n}} \int_{-1}^{C^{2}} \operatorname{up}\left(2^{n+1} \tau+2^{n+1}+1-2 j\right)\left(-1+k 2^{-n}-\tau\right)^{n} \mathrm{~d} \tau= \\
=\frac{2^{C_{n+2}-(n+1)}}{n!} \sum_{j=1}^{k} \delta_{i} \int_{-1}^{1} \operatorname{up}(t)\left(-1+k 2^{-n}-\left(t-2^{n+1}-1+2 j\right) \times\right. \\
\left.\times 2^{-n-1}\right)^{n} \mathrm{dt}=\frac{2^{C_{n+2}^{2}}}{n!} \sum_{j=1}^{k} \delta_{j} \int_{-1}^{1} \operatorname{up}(t)\left(k-j+\frac{1}{2}-\frac{t}{2}\right)^{n} \mathrm{dt}= \\
=\frac{2^{-n^{2}-n-1+C_{n+2}^{2}}}{n!} \sum_{j=1}^{k} \delta_{l} \sum_{l=0}^{\left[\frac{n}{2}\right]} C_{n}^{2 l}\left(k-j+\frac{1}{2}\right)^{n-2 t} a_{2 l^{2}}{ }^{-2 l} . \text { (3.67) } \tag{3.67}
\end{gather*}


Для точек вида $-1+2^{-n}$ поступим следующим образом [59]:


\begin{gather*}
\operatorname{up}\left(-1+2^{-n}\right)=\frac{1}{(n-1)!} \int_{-1}^{-1+2^{-n}} \operatorname{up}^{(n)}(\tau)\left(-1+2^{-n}-\tau\right)^{n-1} \mathrm{~d} \tau= \\
=\frac{2^{c_{n+1}^{2}}}{(n-1)!} \int_{-1}^{-1+2^{-n}} \operatorname{up}\left(2^{n} \tau+2^{n}+1-2\right)\left(-1+2^{-n}-\tau\right)^{n-1} \mathrm{~d} \tau= \\
=\frac{2^{c_{n+1}^{2}-n}}{(n-1)!} \int_{0}^{1} \operatorname{up}(t) t^{n-1} \mathrm{dt} \cdot 2^{-n(n-1)}=\frac{b_{n-1}}{2^{c_{n}^{2}}(n-1)!} . \tag{3.68}
\end{gather*}


По формуле (3.63) найдем\\
$\operatorname{up}^{(l)}\left(-1+2^{-n}\right)=2^{c_{l+1}^{2}} \operatorname{up}\left(-1+2^{l-n}\right)=\frac{2^{c_{l+1}^{2}-c_{n-l}^{2}}}{(n-l-1)!} b_{n-l-1}$.

Получим.теперь формулу для вычисления функции up ( $x$ ) при любом $x$. Достаточно рассмотреть случай $-1 \leqslant x \leqslant 0$, так как функция up $(x)$ четная. Для функции $\varphi_{n}(x)=\operatorname{up}(x)+\operatorname{up}\left(x-2^{-n}\right)$ на интервале $\left[-1+2^{-n},-1+2^{-n+1}\right]$ получаем $\varphi_{n}^{n+1}(x) \equiv 0$, так как $\delta_{2}=-\delta_{1}$. Следовательно, на этом интервале


\begin{equation*}
\varphi_{n}(x)=\sum_{l=0}^{n} \frac{\varphi_{n}^{(l)}\left(-1+2^{-n}\right)}{l l}\left(x+1-2^{-n}\right)^{l}, \tag{3.70}
\end{equation*}


причем $\varphi_{n}^{(l)}\left(-1+2^{-n}\right)=\operatorname{up}^{(l)}\left(-1+2^{-n}\right)$, так как $\operatorname{up}^{(l)}(-1)=0$. Таким образом, для $x \in\left[-1+2^{-n},-1+2^{-n+1}\right]$


\begin{equation*}
\operatorname{up}(x)=\sum_{l=0}^{n} \frac{\operatorname{up}^{(l)}\left(-1+2^{-n}\right)}{l!}\left(x+1-2^{-n}\right)-\operatorname{up}\left(x-2^{-n}\right), \tag{3.71}
\end{equation*}


т. е. для того чтобы вычислить функцию up $(x)$ на $\left[-1+2^{-n}\right.$, $\left.-1+2^{-n+1}\right]$, нужно знать ее значения на $\left[-1,-1+2^{-n}\right]$. Разбивая интервал $\left[-1,-1+2^{-n}\right]$ на интервалы $\left[-1,-1+2^{-n-1}\right]$ и $\left[-1+2^{-n-1},-1+2^{-n}\right]$ и повторяя те же выкладки, можем свести вычисление функции up $(x)$ на интервале $\left[-1+2^{-n-1},-1+\right. +2^{-n+1}$ ] к вычислению ее на $\left[-1,-1+2^{-n-1}\right]$. Тем самым получаем возможность записать некоторый ряд для вычисления значений функции up ( $x$ ), причем ряд быстро сходится. Действительно, на $\left[-1,-1+2^{-n}\right]$

$$
\begin{aligned}
& 0 \leqslant \operatorname{up}(x) \leqslant \operatorname{up}\left(-1+2^{-n}\right), \\
& \operatorname{up}\left(-1+2^{-n}\right)=\frac{b_{n-1}}{2^{c_{n}^{2}}(n-1)!} .
\end{aligned}
$$

Следовательно, использовав оценку (3.38) для $c_{2 k}(a)$, получим


\begin{equation*}
\operatorname{up}\left(-1+2^{-n}\right) \leqslant \frac{a_{2}\left[\frac{n-1}{2}\right]}{2^{c_{n}^{2}}(n-1)!}<\frac{c_{9}}{2^{c_{n}^{2}+2\left[\frac{n-1}{2}\right]\left(2\left[\frac{n-1}{2}\right]+1\right)}} . \tag{3.72}
\end{equation*}


Выпишем этот ряд для $0 \leqslant x<1$ :\\
$\operatorname{up}(x-1)=\sum_{k=1}^{\infty}(-1)^{s_{k}+1} p_{k} \sum_{j=0}^{k} \frac{2^{C_{j+1}^{2}} b_{k-j-1}}{2^{C_{k-j}^{2}}(k-j-1)!j!}\left(x-0, p_{1} \ldots p_{k}\right)^{i}$;


\begin{equation*}
b_{-1}=1, \quad 0!=(-1)!=1, \quad s_{k}=\sum_{j=1}^{k} p_{j} . \tag{3.73}
\end{equation*}


Здесь величины $b_{l}, l>0$, вычисляются по формулам (3.65), (3.66), а $x=0, p_{1} \ldots p_{k} \ldots$ - запись числа $x$ в двоичной системе счисления.

Таким образом, $x-0, p_{1} \ldots p_{k}=x-\left[x 2^{k}\right] 2^{-k}$. Остаточный член


\begin{equation*}
R_{n}(x)=\operatorname{up}(x-1)-\sum_{k=1}^{n}(-1)^{s_{k}+1} \sum_{i=0}^{k} \frac{2^{c_{i+1}^{2}} b_{k-j-1}\left(x-0, p, \ldots p_{k}\right)^{i}}{2^{c_{k-j}^{2}}(k-j-1)!j!} \tag{3.74}
\end{equation*}


не превосходит up ( $-1+2^{-n}$ ), т. е.


\begin{equation*}
\left|R_{n}(x)\right| \leqslant \frac{C}{2^{2}+2\left[\frac{n-1}{2}\right]\left(2^{\left[\frac{n-1}{2}\right]+1}\right)_{(n-1)!}} \tag{3.75}
\end{equation*}


\section*{§ 4. ПРЕДСТАВЛЕНИЕ МНОГОЧЛЕНОВ В ВИДЕ ЛИНЕЙНЫХ КОМБИНАЦИЙ СДВИГОВ ФУНКЦИИ up( $x$ )}
Из формулы (3.63) следует, что для того чтобы линейная комбинация сдвигов функции up $(x)$


\begin{equation*}
\varphi(x)=\sum_{k=-\infty}^{\infty} c_{k} \text { up }\left(x-k 2^{-n}\right) \tag{3.76}
\end{equation*}


была многочленом степени $n$, необходимо и достаточно выполнение системы равенств


\begin{equation*}
\sum_{j=1}^{2^{n+1}} c_{k+j} \delta_{j}=0, \quad k \in \mathbb{Z} \tag{3.77}
\end{equation*}


т. е. коэффициенты $c_{k}$ должны быть решением разностной однородной системы с постоянными коэффициентами (3.77). Поэтому должно выполняться условие


\begin{equation*}
c_{k}=\sum_{i=1}^{N} \sum_{l=0}^{p_{l}-1} c_{i l} k^{l}\left(\omega_{i}^{(n)}\right)^{k}, \tag{3.78}
\end{equation*}


где $\omega_{i}^{(n)}$ - корни уравнения (степени $2^{n+1}-1$ )


\begin{equation*}
Q_{n}(\lambda)=\sum_{k=1}^{2^{n+1}} \delta_{k} \lambda^{2^{n+1}-k}=0 \tag{3.79}
\end{equation*}


а $p_{i}$ - кратность корня $\omega_{i}^{(n)}$. Вычислим корни уравнения (3.79). Очевидно,


\begin{equation*}
Q_{0}(\lambda)=\lambda-1, \quad \omega_{1}^{(0)}=1 . \tag{3.80}
\end{equation*}


Кроме того, легко видеть, что многочлены $Q_{n}(\lambda)$ удовлетворяют рекуррентному соотношению


\begin{equation*}
Q_{n}(\lambda)=Q_{n-1}(\lambda)\left(\lambda^{2^{n}}-1\right) . \tag{3.81}
\end{equation*}


Поэтому корни характеристического уравнения (3.79) равны:


\begin{equation*}
\omega_{j}^{(n)}=\exp \left(2 \pi i j 2^{-l}\right), \tag{3.82}
\end{equation*}


тде $0 \leqslant j \leqslant 2^{l}, l=0,1, \ldots, n$. Окончательно получаем


\begin{equation*}
c_{k}=\sum_{i=1}^{2^{n}}\left[\cos \left(2 \pi j k 2^{-n}\right) \sum_{l=0}^{v_{2}(j)} c_{i l} k^{l}+\sin \left(2 \pi j k 2^{-n}\right) \sum_{l=0}^{v_{2}(j)} d_{j l} k^{l}\right], \tag{3.83}
\end{equation*}


где $v_{2}(j)-2$-адический показатель числа $j$, т. е. наибольшее целое число $v_{2}(j)$ такое, что $2^{v_{2}(j)}$ делит $j$. В частности, если $c_{k}$ - многочлен степени $n$ от индекса $k$, то функция $\sum_{k} c_{k}$ up $\left(x-k 2^{-n}\right)= =P_{n}(x)$ - многочлен степени $n$. Действительно, корень 1 имеет кратность $n+1\left(j=2^{n}\right)$. В работе [53] приведены формулы для представления $x^{n}$ в виде линейной комбинации сдвигов функции up (x).

Докажем следующую теорему о единственности представляющей функции (т. е. функции, линейные комбинации сдвигов которой могут представлять собой многочлен сколь угодно высоких степеней). Эта теорема является усилением теорем, приведенных в работах [53, 66].

Теорема. Пусть $\varphi(x)$ - финитная функция с носителем [-1,1] класса $C^{\infty}$, удовлетворяющая следующим условиям:\\
1)

$$
\int_{-1}^{1} \varphi(x) \mathrm{d} x=1
$$

2.

$$
\frac{\lim _{n \rightarrow \infty}}{\| \varphi^{(n)}{ }^{n} C_{|-1,1|}}=0 ;
$$

\begin{enumerate}
  \setcounter{enumi}{2}
  \item для каждого натурального $n$ существуют коэффициенты $c_{j}^{(n)}$ такие, что
\end{enumerate}

$$
\sum_{i=-\infty}^{\infty} c_{i}^{(n)} \varphi\left(x-j 2^{-n}\right)=x^{n} .
$$

Тогда $\varphi(x) \equiv \operatorname{up}(x)$.\\
Доказательство. Из условия 1 теоремы следует

$$
\sum_{i=-\infty}^{\infty} c_{i}^{(0)} \varphi(x-j) \equiv 1,
$$

где коэффициенты $c_{j}^{(0)}$ обязательно равны 1. Отсюда $\varphi^{\prime}(x)= =2 \varphi_{1}(2 x+1)-2 \varphi_{1}(2 x-1)$, где $\varphi_{1}(x)$ - финитная функция с носителем $[-1,1]$ класса $C^{\infty}$, удовлетворяющая всем условиям теоремы. Аналогично находим

$$
\varphi_{1}^{\prime}(x)=2 \varphi_{2}(2 x+1)-2 \varphi_{2}(2 x-1),
$$

где $\varphi_{2}(x)$ удовлетворяет тем же условиям, и т. д. Получаем последовательность функций $\left\{\varphi_{n}(x)\right\}$, удовлетворяющих условиям теоремы и таких, что


\begin{equation*}
\varphi_{n}^{\prime}(x)=2 \varphi_{n+1}(2 x+1)-2 \varphi_{n+1}(2 x-1) . \tag{3.84}
\end{equation*}


Проинтегрируем (3.84):


\begin{equation*}
\varphi_{n}(x)=2 \int_{-1}^{x}\left(\varphi_{n+1}(2 t+1)-\varphi_{n+1}(2 t-1)\right) d t . \tag{8.85}
\end{equation*}


Из определения функции up $(x)$ имеем


\begin{equation*}
\operatorname{up}(x)=2 \int_{-1}^{x}(\operatorname{up}(2 t+1)-\operatorname{up}(2 t-1)) \mathrm{dt} \tag{3.86}
\end{equation*}


Из условия 2 теоремы следует


\begin{equation*}
\frac{\lim _{n \rightarrow \infty}}{n \rightarrow \varphi_{n+1} \|_{C}}=_{2^{n}} \frac{2^{-C_{n+2}^{2}}\left\|\varphi^{(n+1)}(x)\right\|_{C}}{2^{n}}=0 . \tag{3.87}
\end{equation*}


Кроме того, $\varphi_{n}(0)=1, \varphi_{n}(x-1)+\varphi_{n}(x)=1$ при $x \in[0,1]$. То же верно и для функции up $(x)$. Положим


\begin{equation*}
\gamma_{n}(x)=\operatorname{up}(x)-\varphi_{n}^{-}(x) . \tag{3.88}
\end{equation*}


Согласно (3.87)


\begin{equation*}
\varliminf_{n \rightarrow \infty} \frac{\left\|\gamma_{n+1}(x)\right\|_{C}}{2^{n}}=0 . \tag{3.89}
\end{equation*}


Функция $\gamma_{n}(x)$ обращается в 0 в точках $-1,0,1$ и, кроме того, для $x \in[0,1]$


\begin{equation*}
\gamma_{n}(x)=-\gamma_{n}(x-1) . \tag{3.90}
\end{equation*}


Из (3.84) и (3.88) находим


\begin{equation*}
\gamma_{n}(x)=2 \int_{-1}^{x}\left(\gamma_{n+1}(2 t+1)-\gamma_{n+1}(2 t-1)\right) \mathrm{dt} . \tag{3.91}
\end{equation*}


При $x \in\left[-1,-\frac{1}{2}\right]$ из (3.91) следует


\begin{equation*}
\left|\gamma_{n}(x)\right| \leqslant\left\|\gamma_{n+1}\right\|_{c} . \tag{3:92}
\end{equation*}


Согласно (3.90) неравенство (3.92) выполняется и на $\left[-\frac{1}{2}, 0\right]$ и, следовательно, на $[-1,1]$. Кроме того, в силу (3.90) на $[-1,0]$


\begin{equation*}
\left|\gamma_{n}(x)\right| \leqslant\left\|\gamma_{n+1}\right\| c \sigma_{1}(2 x+1), \tag{3.93}
\end{equation*}


где $\sigma_{1}(x)=1-|x|$. Поэтому


\begin{equation*}
\left\|\gamma_{n-1}(x)\right\|_{c} \leqslant \frac{1}{2}\left\|\varphi_{n+1}\right\|_{c} \tag{3.94}
\end{equation*}


и


\begin{equation*}
\left|\gamma_{n-1}(x)\right| \leqslant \frac{1}{2}\left\|\gamma_{n+1}\right\| c \sigma_{2}(2 x+H), \tag{3.95}
\end{equation*}


где


\begin{equation*}
\sigma_{n}(x)=2 \int_{-1}^{x}\left(\sigma_{n-1}(2 t+1)-\sigma_{n-1}(2 t-1)\right) \mathrm{dt} . \tag{3.96}
\end{equation*}


Продолжая эти рассуждения, получаем


\begin{equation*}
\left\|\gamma_{0}\right\|_{c}=\|\operatorname{up}(x)-\varphi(x)\|_{c} \leqslant\left(\frac{1}{2}\right)^{n}\left\|\gamma_{n+1}\right\| c . \tag{3.97}
\end{equation*}


Здесь использованы условия $0 \leqslant \sigma_{n}(x) \leqslant 1 ; \int_{-1}^{0} \sigma_{n}(x) \mathrm{dx}=\frac{1}{2}$.\\
Поэтому из условия 3 теоремы и (3.89) следует $\left\|\gamma_{0}\right\|_{c}=0$, т. е. $\varphi(x) \equiv \mathrm{up}(x)$. Теорема доказана.

Замечание. Ограничение на рост производных в условии 3 теоремы нельзя ослабить. Например, функция $\varphi(x)=\operatorname{up}(x)+\alpha \mathrm{up}^{\prime}(x)$, $\alpha \neq 0$, удовлетворяет условиям 1 и 3 теоремы, но

$$
\lim _{n \rightarrow \infty} \frac{\left\|\varphi^{(n)}(x)\right\|_{C}}{2^{C_{n+2}^{2}}}=|\alpha|
$$

и, конечно, теорема для нее не верна: up $(x) \neq u$ p $(x)+\alpha$ up $^{\prime}(x)$. Функции $\sigma_{n}(x)$, используемые при доказательстве теоремы,- финитные сплайны - подробно изучались в работах [62, 63].

\section*{§ 5. ФУнкции $\operatorname{Fup}_{n}(x)$}
Рассмотрим преобразование Фурье функции


\begin{align*}
F(t)=F_{2}(t)= & \prod_{j=1}^{\infty} \frac{\sin t 2^{-j}}{t 2^{-j}}=\frac{2 \cos \frac{t}{4} \sin \frac{t}{4}}{\frac{t}{2}} \prod_{j=2}^{\infty} \frac{\sin t 2^{-1}}{t 2^{-j}}= \\
& =\cos \frac{t}{4}\left(\frac{\sin \frac{t}{4}}{\frac{t}{4}}\right) \prod_{j=3}^{\infty} \frac{\sin t 2^{-j}}{t 2^{-j}} . \tag{3.98}
\end{align*}


Положим


\begin{align*}
& F^{[1]}(t)=\left(\frac{\sin \frac{t}{4}}{\frac{t}{4}}\right)^{2} \prod_{j=3}^{\infty} \frac{\sin t 2^{-j}}{t 2^{-j}} ;  \tag{3.99}\\
& F \operatorname{up}_{1}(x)=\frac{1}{2 \pi} \int_{-\infty}^{\infty} e^{-i t x} F^{[1]}(t) \mathrm{dt} .
\end{align*}


Тогда из (3.98) следует


\begin{equation*}
\operatorname{up}(x)=\left(F \operatorname{up}_{1}\left(x-\frac{1}{4}\right)+F \operatorname{up}_{1}\left(x+\frac{1}{4}\right)\right) 2^{-1} . \tag{3.100}
\end{equation*}



\end{document}